\chapter{Conclusion}

\section{Conclusion}
This thesis presented the design, implementation, and empirical evaluation of a real-time Camera AI Attendance System using face recognition. The project demonstrates the feasibility of combining computer vision models with conventional software systems to automate check-in and check-out workflows.

Key achievements of the project include:
\begin{itemize}
    \item \textbf{Efficient Multi-Model Pipeline:} We integrated YOLOv8-face and InceptionResnetV1 to achieve reliable face verification. Using an ROI tracking fallback (\texttt{cv2.matchTemplate}) reduced CNN detection overhead by 90\%, boosting system performance from 13.8 FPS to 25.2 FPS on standard CPU hardware.
    \item \textbf{Lightweight Emotion-Gated Interactions:} A custom \textit{Mini-Xception} model was trained on merged public datasets (FER-2013, FERPlus, CK+, and JAFFE). Calibration techniques (Temperature Scaling and Prior Weight Correction) reduced model bias and improved the F1-scores of rare emotion classes by up to 15\%. This enabled a smiling check-in gate.
    \item \textbf{Rule-Based Liveness and Gesture Filters:} A rule-based 2D anti-spoofing filter was implemented, combining grayscale contrast, Laplacian frequency bandpass, skin-color ratio, glare, and bezel checks to reject printed or digital photos. A hand-waving gesture check-out module detected motion swings to separate checking-out employees from passers-by.
    \item \textbf{Production-Ready Integration:} The pipeline was integrated with a FastAPI backend and a web dashboard. The dashboard's split sidebar displays check-in (green) and check-out (orange) events in real-time, providing an intuitive interface for administrators.
\end{itemize}

\section{Limitations}
Despite its performance, the system has several limitations:
\begin{itemize}
    \item \textbf{Environmental Sensitivity:} Under extremely dim lighting, strong backlight (e.g., facing a window), or severe facial angles ($>45^\circ$), face detection recall and verification accuracy degrade.
    \item \textbf{2D Spoofing Vulnerabilities:} The rule-based 2D liveness filter detects simple screen replay and paper print attacks. However, it remains vulnerable to high-fidelity 3D masks, silicon face overlays, or video playback on large screens without visible bezels, where Canny edge detection cannot locate borders.
    \item \textbf{Database Scalability:} For small to medium organizations ($<500$ employees), the database search using cosine similarity is highly efficient. However, scaling to larger corporate databases with thousands of employees would require indexing optimization (such as FAISS) to maintain low latency.
    \item \textbf{CPU Overload under Multi-Streams:} Although a single camera runs smoothly at 25+ FPS, running multiple concurrent camera streams on the same CPU would saturate resources, leading to processing delays and frame drops.
\end{itemize}

\section{Future Work}
To address these limitations, future work will focus on:
\begin{itemize}
    \item \textbf{Depth-Based Liveness Detection:} Integrating structured light or time-of-flight (ToF) 3D cameras to perform depth map analysis. This would prevent physical 2D spoofing attacks.
    \item \textbf{Edge Hardware Acceleration:} Deploying the model pipeline on dedicated edge accelerators, such as NVIDIA Jetson Nano or Google Coral Edge TPU, to run heavy neural network models (e.g., InceptionResnetV1) at 30+ FPS with lower CPU usage.
    \item \textbf{Multi-Camera Synchronization:} Extending the backend to support distributed camera streams, allowing employees to check in at one entrance and check out at another with synchronized database records.
    \item \textbf{Advanced Spoofing Models:} Incorporating lightweight deep learning-based liveness networks (e.g., FaceAntiSpoofing) to complement the rule-based filters.
\end{itemize}
